% Part 1 of Day 15: Course summary.
%
% Topics of the syllabus: review of the main trade-offs from the three weeks;
% questions. Parts 2 and 3 of the lecture time are the capstone build: teams
% work, and the instructor visits each team. The syllabus gives no exercise
% for Part 1: the exercise prepares the trade-off section of the capstone
% report.
%
% This part repeats no frame of an earlier day. It puts the results of the
% days side by side. Each value comes from an experiment of the course or
% from a published source that the day named. The plan of the course
% repository lists each value with its experiment script. The source line of
% each frame names the days.

\theorypart{Course summary}%
  {name one capstone trade-off and the number that decides it}

% ------------------------------------------------------------------ map
\begin{frame}{Fifteen days in one view}
  \begingroup\centering
  \begin{tikzpicture}[font=\scriptsize,
      wk/.style={draw=coolgray, rounded corners=3pt, fill=boxgray,
                 align=left, text width=3.0cm, minimum height=2.7cm,
                 anchor=north}]
    \node[wk, draw=cardinalred] at (0,0) {%
      \textbf{Week 1: TinyML on\\ the XIAO}\\[3pt]
      1 Landscape, budgets\\ 2 Sensors, data\\ 3 Model to the board\\
      4 Quantization\\ 5 Audio and vision};
    \node[wk, draw=treegreen] at (3.6,0) {%
      \textbf{Week 2: optimization,\\ Raspberry Pi}\\[3pt]
      6 Pruning, distillation\\ 7 Runtimes, graphs\\ 8 Object detection\\
      9 Benchmarking, energy\\ 10 Language models};
    \node[wk, draw=darkcyan] at (7.2,0) {%
      \textbf{Week 3: connected\\ systems}\\[3pt]
      11 Networks, RTSP\\ 12 Decisions, MQTT\\ 13 Monitoring\\
      14 Production design\\ 15 Capstone};
  \end{tikzpicture}\par\endgroup
  \vskip 0.5em
  \small
  \begin{itemize}
    \item Week 1 put one model on one microcontroller. Week 2 made models
          smaller and faster, and measured them. Week 3 joined devices into
          one system that runs for a long time.
    \item Each day measured at least one trade-off. This part puts the
          measured numbers side by side.
  \end{itemize}
\end{frame}

% ------------------------------------------------------------------ workflow
\begin{frame}{The workflow of Day 1, with its days}
  \begingroup\centering
  \begin{tikzpicture}[font=\scriptsize, >=stealth,
      st/.style={draw=coolgray, rounded corners=2pt, fill=boxgray,
                 align=center, text width=1.15cm, minimum height=0.8cm}]
    \foreach \i/\n/\d in {0/Collect data/{2, 5, 8}, 1/Train/{3, 6, 8},
        2/Optimize/{4, 6}, 3/Convert/{3, 7}, 4/Deploy/{3, 7, 10, 11, 12},
        5/Monitor/{13, 14}} {
      \node[st] (s\i) at (1.75*\i,0) {\n};
      \node[text=cardinalred, align=center, text width=1.45cm] at
        (1.75*\i,-0.9) {Days \d};
    }
    \foreach \i [evaluate=\i as \j using int(\i+1)] in {0,...,4} {
      \draw[->] (s\i) -- (s\j);
    }
    \draw[->, darkcyan] (s5.north) -- ++(0,0.45) -| node[pos=0.25, above]
      {new data, new model (Days 13, 14)} (s0.north);
  \end{tikzpicture}\par\endgroup
  \vskip 0.4em
  \small
  \begin{itemize}
    \item Day 9 measured the stages: latency, memory, energy, accuracy.
    \item Day 1 said: the device constraints come first. Each later day
          showed one more constraint: RAM, flash, latency, energy, the
          network, the time of the people who run the system.
    \item The loop does not end at the deployment. A deployed model needs
          monitoring and updates.
  \end{itemize}
\end{frame}

% ------------------------------------------------------------------ list
\begin{frame}{The main trade-offs}
  \footnotesize
  \renewcommand{\arraystretch}{1.15}
  \setlength{\tabcolsep}{4pt}
  \begin{tabular}{@{}L{0.32\textwidth}L{0.30\textwidth}L{0.28\textwidth}@{}}
    \toprule
    \textbf{Trade-off} & \textbf{You gain} & \textbf{You pay} \\
    \midrule
    Compute on the device & Latency, privacy, data volume & Memory,
      energy, a smaller model \\
    Fewer bits, fewer weights & Size, often speed & Accuracy, a
      calibration step \\
    Smaller input & Speed & Small objects \\
    Local decision rules & Fewer false actions & Delay \\
    Reliable delivery & No lost message & Bytes, double messages \\
    Monitoring and logs & Early warning & Storage, false alerts \\
    Security and updates & Safe operation & Bytes, time, work \\
    \bottomrule
  \end{tabular}
  \vskip 0.6em
  \small
  No row has a free choice. The requirements of the product decide which
  side wins. The next frames give the measured numbers of each row.
\end{frame}

% ------------------------------------------------------------------ where
\begin{frame}{Where to compute: the data that leaves}
  \begingroup\centering
  \begin{tikzpicture}[font=\scriptsize, x=0.75cm, y=0.55cm]
    % Bar length = log10(bytes in one day).
    \foreach \x in {3,5,7,9} {
      \draw[coolgray!50] (\x,-0.4) -- (\x,3.0);
      \node[below, coolgray] at (\x,-0.4) {$10^{\x}$};
    }
    \fill[cardinalred!60] (0,2.1) rectangle (9.442,2.7);
    \fill[darkcyan!60]    (0,1.1) rectangle (7.320,1.7);
    \fill[treegreen!60]   (0,0.1) rectangle (3.782,0.7);
    \node[anchor=west, text=white] at (0.1,2.4) {Raw audio, 16 kHz: 2\,764\,800\,000};
    \node[anchor=west, text=white] at (0.1,1.4) {A result each 0.5 s: 20\,908\,800};
    \node[anchor=west, fill=white, inner sep=1pt] at (3.9,0.4) {50 keyword events: 6050};
  \end{tikzpicture}\par\endgroup
  \vskip 0.2em
  \small
  \begin{itemize}
    \item Bytes in one day of one microphone (log scale). Local inference
          sends results, not raw data: a factor of 132 for the results and
          of about 457\,000 for the events.
    \item The radio of the ESP32-S3 needs up to 1122 mW to send; the
          processor needs 217 mW. Each byte that stays costs no radio
          energy.
    \item The price: the model must fit the device (next frame).
  \end{itemize}
  \source{Day 12, Part 3 (data volumes); Day 9 (ESP32-S3 data sheet 2.2,
    the power is a calculation).}
\end{frame}

% ------------------------------------------------------------------ budgets
\begin{frame}{The three budgets of a model}
  \footnotesize
  \renewcommand{\arraystretch}{1.15}
  \setlength{\tabcolsep}{4pt}
  \begin{tabular}{@{}L{0.30\textwidth}rrr@{}}
    \toprule
    \textbf{Model} & \textbf{Parameters} & \textbf{MACs} &
      \textbf{Peak activations} \\
    \midrule
    ResNet-18 & 11\,689\,512 & 1\,814\,073\,344 & 1\,003\,520 \\
    MobileNetV2 & 3\,504\,872 & 300\,774\,272 & 1\,505\,280 \\
    Small depthwise CNN & 14\,186 & 5\,116\,160 & 110\,592 \\
    \bottomrule
  \end{tabular}
  \vskip 0.6em
  \small
  \begin{itemize}
    \item The parameters fill the flash, the operations take the time and
          the energy, and the peak activations fill the RAM.
    \item The budgets are independent: MobileNetV2 has 3.3 times fewer
          parameters than ResNet-18, but a 1.5 times larger peak.
    \item Of the three, only the small depthwise CNN in \code{int8} fits a
          microcontroller with 256 KB of RAM.
  \end{itemize}
  \keypoint{Check all three budgets before you train. A model that misses
    one budget does not run.}
  \source{Day 1, lab notebook (calculated, values for one image).}
\end{frame}

% ------------------------------------------------------------------ bits
\begin{frame}{Fewer bits: size against accuracy}
  \footnotesize
  \renewcommand{\arraystretch}{1.15}
  \setlength{\tabcolsep}{4pt}
  \begin{tabular}{@{}L{0.40\textwidth}L{0.24\textwidth}L{0.24\textwidth}@{}}
    \toprule
    \textbf{Model} & \textbf{\code{float32}} & \textbf{\code{int8}} \\
    \midrule
    Motion model of Day 3 (file) & 26\,732 bytes & 8680 bytes \\
    MobileNetV2 (accuracy, Imagenette) & 77.4 percent & 75.2 percent \\
    YOLO11n (mAP50-95, 500 COCO images) & 0.424 & 0.331 \\
    Raw window model of Day 3, calibration with one class & 100.0 percent
      & 69.2 percent \\
    \bottomrule
  \end{tabular}
  \vskip 0.6em
  \small
  \begin{itemize}
    \item \code{int8} makes the files 3 to 3.5 times smaller. The loss of
          accuracy depends on the model: small for a classifier, large for
          a detector.
    \item The calibration data decides the ranges. Data of one class gave
          a large loss.
    \item Quantization-aware training (Day 4) and \code{int8} weights with
          \code{float32} activations reduce the loss.
  \end{itemize}
  \source{Day 4 (motion model, calibration); Day 7 (MobileNetV2); Day 8
    (YOLO11n). Work computer, LiteRT.}
\end{frame}

% ------------------------------------------------------------------ smaller
\begin{frame}{Fewer weights: what saves what}
  \footnotesize
  \renewcommand{\arraystretch}{1.15}
  \setlength{\tabcolsep}{4pt}
  \begin{tabular}{@{}L{0.46\textwidth}rrr@{}}
    \toprule
    \textbf{Fashion-MNIST CNN (\code{float32})} & \textbf{File} &
      \textbf{Accuracy} & \textbf{Time} \\
    \midrule
    Baseline, filters 16-32-64 & 99\,320 B & 89.1\,\% & 71\,$\mu$s \\
    90 percent of the weights set to zero & 99\,320 B & 76.4\,\% & no
      gain \\
    Filters removed: 8-16-32 & 28\,408 B & 85.6\,\% & 26\,$\mu$s \\
    Separable layers & 20\,872 B & 86.3\,\% & 27\,$\mu$s \\
    \bottomrule
  \end{tabular}
  \vskip 0.6em
  \small
  \begin{itemize}
    \item Zeros in a dense file save nothing on a normal processor. Remove
          whole filters, or design the model small.
    \item A student trained only on the outputs of a teacher (90.3
          percent) reached 85.3 percent: the student needs no labels, and a
          copy of a model needs only its outputs (Day 14).
  \end{itemize}
  \source{Day 6, Parts 1 to 3 (work computer, LiteRT, one thread; a time
    changes by 5 to 10 $\mu$s between two runs).}
\end{frame}

% ------------------------------------------------------------------ runtime
\begin{frame}{The runtime matters as much as the model}
  \footnotesize
  \renewcommand{\arraystretch}{1.15}
  \setlength{\tabcolsep}{4pt}
  \begin{tabular}{@{}L{0.55\textwidth}r@{}}
    \toprule
    \textbf{MobileNetV2, one thread, the same weights} & \textbf{Latency} \\
    \midrule
    LiteRT with its reference kernels & 362.7 ms \\
    LiteRT with its built-in kernels & 41.8 ms \\
    PyTorch & 31.9 ms \\
    ONNX Runtime with no graph optimization & 26.4 ms \\
    NCNN & 16.6 ms \\
    ONNX Runtime & 11.7 ms \\
    LiteRT with XNNPACK & 11.7 ms \\
    \bottomrule
  \end{tabular}
  \vskip 0.6em
  \small
  \begin{itemize}
    \item A factor of 31 between the slowest and the fastest runtime, with
          the same model.
    \item Fused operators, a good memory layout, and SIMD kernels make the
          difference (Day 7, Part 3).
  \end{itemize}
  \source{Day 7, Part 2 (x86 processor with AVX2, median of 50 runs). The
    values say nothing about a board.}
\end{frame}

% ------------------------------------------------------------------ resolution
\begin{frame}{A smaller input: speed against small objects}
  \begin{columns}[c]
    \begin{column}{0.56\textwidth}
      \centering
      \begin{tikzpicture}[font=\scriptsize, x=0.62cm, y=4.2cm]
        \draw[->, coolgray] (0,0) -- (7.2,0);
        \node[coolgray] at (3.6,-0.17) {GFLOP of one image};
        \draw[->, coolgray] (0,0) -- (0,0.88);
        \foreach \y in {0.2,0.4,0.6,0.8} {
          \draw[coolgray] (-0.08,\y) -- (0,\y) node[left] {\y};
        }
        \foreach \x in {2,4,6} {
          \draw[coolgray] (\x,-0.02) -- (\x,0) node[below=1pt] {\x};
        }
        \draw[treegreen, thick] (0.80,0.62) -- (1.62,0.72) -- (3.67,0.77) --
          (6.54,0.79);
        \draw[darkcyan, thick] (0.80,0.244) -- (1.62,0.319) -- (3.67,0.391) --
          (6.54,0.417);
        \draw[cardinalred, thick] (0.80,0.01) -- (1.62,0.07) -- (3.67,0.16) --
          (6.54,0.23);
        \foreach \x/\y in {0.80/0.62, 1.62/0.72, 3.67/0.77, 6.54/0.79,
            0.80/0.244, 1.62/0.319, 3.67/0.391, 6.54/0.417, 0.80/0.01,
            1.62/0.07, 3.67/0.16, 6.54/0.23} {
          \fill (\x,\y) circle (1.2pt);
        }
        \node[treegreen, anchor=north east] at (6.6,0.76) {recall, large};
        \node[darkcyan, anchor=north] at (5.1,0.385) {mAP50-95};
        \node[cardinalred, anchor=north east] at (6.6,0.205) {recall,
          small};
        \foreach \x/\y/\l in {0.80/0.244/224, 1.62/0.319/320, 3.67/0.391/480,
            6.54/0.417/640} {
          \node[coolgray, font=\tiny, above=1pt] at (\x,\y) {\l};
        }
      \end{tikzpicture}
    \end{column}
    \begin{column}{0.40\textwidth}
      \small
      \begin{itemize}
        \item YOLO11n at 640, 480, 320, and 224 pixels (grey numbers).
        \item From 640 to 320 pixels: 4 times fewer operations; the recall
              of large objects falls from 0.79 to 0.72, of small objects
              from 0.23 to 0.07.
        \item Select the size from the size of the objects in the image,
              not from the frame rate only.
      \end{itemize}
    \end{column}
  \end{columns}
  \source{Day 8, Parts 2 and 3 (500 COCO validation images; recall at the
    score threshold 0.25).}
\end{frame}

% ------------------------------------------------------------------ measurement
\begin{frame}{How you measure changes the number}
  \footnotesize
  \renewcommand{\arraystretch}{1.15}
  \setlength{\tabcolsep}{4pt}
  \begin{tabular}{@{}L{0.58\textwidth}r@{}}
    \toprule
    \textbf{MobileNetV2, LiteRT, one thread} & \textbf{Value} \\
    \midrule
    Mean of the first 5 inferences of a new process & 13.35 ms \\
    Median of 2000 inferences after a warm-up & 11.621 ms \\
    p99 alone on the core & 12.264 ms \\
    p99 with a second program on the core & 27.452 ms \\
    \midrule
    95 percent interval of an accuracy, 100 test samples & 8.2 points \\
    The same, 1000 test samples & 2.6 points \\
    \bottomrule
  \end{tabular}
  \vskip 0.6em
  \small
  \begin{itemize}
    \item Warm up, repeat, and give a distribution: the median and a high
          percentile.
    \item A second program changes the tail, not the median.
    \item A small test set cannot show a difference of a few points.
  \end{itemize}
  \source{Day 9, Parts 1 and 2 (x86 processor with AVX2).}
\end{frame}

% ------------------------------------------------------------------ energy
\begin{frame}{Energy: the rest state decides}
  \footnotesize
  \renewcommand{\arraystretch}{1.15}
  \setlength{\tabcolsep}{4pt}
  \begin{tabular}{@{}L{0.40\textwidth}rrrr@{}}
    \toprule
    \textbf{Duty cycle (light sleep between)} & \textbf{100\,\%} &
      \textbf{10\,\%} & \textbf{1\,\%} & \textbf{0.1\,\%} \\
    \midrule
    Mean power of the chip & 217 mW & 22.4 mW & 2.95 mW & 1.01 mW \\
    Life on 1000 mAh at 3.7 V & 17 h & 6.9 d & 52 d & 153 d \\
    \bottomrule
  \end{tabular}
  \vskip 0.6em
  \small
  \begin{itemize}
    \item A faster model saves energy only when the device sleeps after it.
    \item The board adds its own rest power: 19 mW in light sleep for the
          XIAO ESP32S3 Sense, 24 times the chip. With Wi-Fi connected, at
          least 167.2 mW (Day 14).
    \item A Raspberry Pi 5 needs 3.0 to 3.5 W with no load: a mains
          supply.
  \end{itemize}
  \source{Day 9, Part 3 (ESP32-S3 data sheet 2.2, Seeed board values,
    \emph{Edge AI Engineering}; the life is a calculation).}
\end{frame}

% ------------------------------------------------------------------ llm
\begin{frame}{Language models: memory sets the speed}
  \footnotesize
  \renewcommand{\arraystretch}{1.15}
  \setlength{\tabcolsep}{4pt}
  \begin{tabular}{@{}L{0.44\textwidth}rr@{}}
    \toprule
    & \textbf{\code{llama3.2:1b}} & \textbf{\code{llama3.2:3b}} \\
    \midrule
    Format, file & \code{Q8\_0}, 1.32 GB & \code{Q4\_K\_M}, 2.02 GB \\
    Generation on a Raspberry Pi 5 & 8.99 tokens/s & 5.3 tokens/s \\
    20 commands to JSON with a schema, correct & 18 & 20 \\
    \bottomrule
  \end{tabular}
  \vskip 0.6em
  \small
  \begin{itemize}
    \item Each token reads all weights: the file size times the rate is
          about the memory bandwidth. Fewer bits give more tokens each
          second.
    \item The larger model was correct more often. The context costs RAM:
          the cache of the 3B model needs 3.76 GB at 32\,768 tokens.
    \item A schema and a tool call fix the form of the output, not the
          facts.
  \end{itemize}
  \source{Day 10 (rates: textbook kit lab; JSON test: Ollama on the work
    computer; sizes: GGUF headers). 1 GB = $10^9$ bytes.}
\end{frame}

% ------------------------------------------------------------------ streams
\begin{frame}{Streams: bit rate, start time, and age}
  \footnotesize
  \renewcommand{\arraystretch}{1.15}
  \setlength{\tabcolsep}{4pt}
  \begin{tabular}{@{}L{0.55\textwidth}L{0.33\textwidth}@{}}
    \toprule
    \textbf{Choice} & \textbf{Measured} \\
    \midrule
    MJPEG against H.264, 1280 x 720 & 11.09 against 2.62 Mbit/s \\
    Keyframe each 1 s against each 10 s: wait of a new reader & median
      0.50 s against 5.51 s \\
    Reader in order against newest frame, model of 100 ms & age 6711 ms
      against 56 ms \\
    Default decoder threads against 1 thread & 538 ms against 40 ms \\
    \bottomrule
  \end{tabular}
  \vskip 0.6em
  \small
  \begin{itemize}
    \item Compression saves bandwidth and costs start time and CPU time.
    \item A slow model must drop frames. A queue in order makes each
          decision older and older.
    \item A default setting can hide a large delay: measure the age of the
          frame, not only the frame rate.
  \end{itemize}
  \source{Day 11, Parts 2 and 3 (work computer, MediaMTX, FFmpeg, OpenCV).}
\end{frame}

% ------------------------------------------------------------------ decisions
\begin{frame}{Decisions and delivery}
  \footnotesize
  \renewcommand{\arraystretch}{1.15}
  \setlength{\tabcolsep}{4pt}
  \begin{tabular}{@{}L{0.40\textwidth}rrr@{}}
    \toprule
    \textbf{Rule on the scores, 60 s} & \textbf{Switch-ons} & \textbf{False}
      & \textbf{Delay on} \\
    \midrule
    One window above 0.5 & 16.11 & 8.60 & 0.04 s \\
    3 of 5 windows & 1.07 & 0.00 & 0.21 s \\
    \bottomrule
  \end{tabular}
  \vskip 0.5em
  \begin{tabular}{@{}L{0.40\textwidth}rrr@{}}
    \toprule
    \textbf{MQTT, a link lost for 1.5 s} & \textbf{QoS 0} & \textbf{QoS 1} &
      \textbf{QoS 2} \\
    \midrule
    Bytes of one IMU message & 147 & 193 & 281 \\
    Lost messages & 51 & 0 & 56 \\
    Double messages & 0 & 20 & 0 \\
    \bottomrule
  \end{tabular}
  \vskip 0.5em
  \small
  \begin{itemize}
    \item A rule over several windows removed the false actions for 0.2 s
          of delay.
    \item QoS 1 and a sequence number to drop doubles gave no loss. QoS 2
          lost messages with this client after a reconnection.
  \end{itemize}
  \source{Day 12, Parts 1 and 2 (simulation, mean of 100 seeds; Mosquitto
    2.0.11, paho-mqtt 2.1.0).}
\end{frame}

% ------------------------------------------------------------------ operations
\begin{frame}{Monitoring and production}
  \footnotesize
  \renewcommand{\arraystretch}{1.15}
  \setlength{\tabcolsep}{4pt}
  \begin{tabular}{@{}L{0.52\textwidth}L{0.36\textwidth}@{}}
    \toprule
    \textbf{Choice} & \textbf{Measured} \\
    \midrule
    Alert on the mean of 10 s against 60 s (5 hours) & 4 against 1 false
      alerts \\
    Log each frame against a summary each 60 s (1 hour) & 7.89 MB against
      18.9 KB \\
    Update of a 9.9 MB folder with restart against resume (100 devices,
      7 days) & 17 against 0 devices with no file \\
    MQTT with TLS against plain: connect & 51.9 ms against 2.4 ms \\
    \bottomrule
  \end{tabular}
  \vskip 0.6em
  \small
  \begin{itemize}
    \item A fast alert is a noisy alert. A longer window gives fewer false
          alerts and a later alert.
    \item A summary with each error keeps the important records for a small
          part of the storage.
    \item Each defense costs bytes or time, but much less than the failure
          that it prevents.
  \end{itemize}
  \source{Day 13, Parts 2 and 3 (simulations); Day 14, Parts 1 and 2
    (simulation, Mosquitto on the work computer).}
\end{frame}

% ------------------------------------------------------------------ patterns
\begin{frame}{Patterns that repeat}
  \small
  \begin{enumerate}
    \item \textbf{Send results, not raw data.} Days 1, 5, 11, 12, 14.
    \item \textbf{Measure on the target.} On the x86 work computer,
          \code{int8} kernels were 1.6 times slower than \code{float32}
          (Day 6). On a Raspberry Pi 5, a published \code{int8} model was
          2.9 times faster (Day 7).
    \item \textbf{A published number is not your number.} Its board, its
          settings, and its window are different (Day 9).
    \item \textbf{Check the defaults.} Decoder threads (Day 11), the
          context length (Day 10), QoS 2 after a reconnection (Day 12),
          the calibration data (Day 4).
    \item \textbf{Test the failure.} Cut the link, darken the image, stop
          the program, fill the disk (Days 12 to 14).
    \item \textbf{Predict, then measure, then explain the difference.} The
          method of each lab.
  \end{enumerate}
  \source{Day 7, Raspberry Pi 5, XNNPACK: 10.84 ms and 3.69 ms.\\
    \emph{EdgeML with Raspberry Pi}, guide for ExecuTorch.}
\end{frame}

% ------------------------------------------------------------------ decision log
\begin{frame}{The Decision Log: one decision, one number}
  \begin{block}{Form (about 100 words)}
    \small
    The decision. The options. The measured numbers of each option. The
    trade-off: what you gain and what you pay. The condition that would
    change the decision.
  \end{block}
  \vskip 0.4em
  \begin{exampleblock}{Example from Day 14, Part 3}
    \small
    We use a window of 2 s. The window is 2.0 s of the alarm budget of
    2.6 s, against the requirement of 3.0 s. A window of 1 s would leave
    more margin, but we did not measure its recall. We keep 2 s until a
    test on the recordings of the site shows the recall of 1 s.
  \end{exampleblock}
  \vskip 0.4em
  \small
  Your capstone report has a benchmark table and the Decision Log of your
  main decisions. The demonstration defends them.
  \source{Method: \emph{Machine Learning Systems}, instructor guide
    (assessment). Budget values of the Day 14 example.}
\end{frame}

% ------------------------------------------------------------------ next
\begin{frame}{After the course}
  \footnotesize
  \renewcommand{\arraystretch}{1.15}
  \setlength{\tabcolsep}{4pt}
  \begin{tabular}{@{}L{0.42\textwidth}L{0.48\textwidth}@{}}
    \toprule
    \textbf{Topic} & \textbf{Backup modules of the course} \\
    \midrule
    Low power and sleep modes & MC-9 \\
    Neural processing units & MC-11 (Grove Vision AI V2), PI-5 \\
    Other runtimes & PI-2 (ExecuTorch), GA-4 (llama.cpp), GA-5 \\
    Learning on the device & SY-4 (on-device and federated learning) \\
    Security in practice & SY-2 (MQTT with TLS), SY-5, SY-7 (OTA) \\
    Generative applications & GA-1 (retrieval), GA-3 (agents), GA-7
      (voice) \\
    Responsible edge AI & SY-6, SY-10 (faults, drift, attacks) \\
    \bottomrule
  \end{tabular}
  \vskip 0.6em
  \small
  \begin{itemize}
    \item Each module has its own deck and lab. Ask the instructor for the
          material.
    \item The textbook (mlsysbook.ai) and the kit labs of the companion
          books continue each topic.
  \end{itemize}
\end{frame}

% ------------------------------------------------------------------ today
\begin{frame}{Today: the capstone build}
  \begin{columns}[T]
    \begin{column}{0.50\textwidth}
      \footnotesize
      \renewcommand{\arraystretch}{1.15}
      \setlength{\tabcolsep}{3pt}
      \begin{tabular}{@{}L{1.5cm}L{3.6cm}@{}}
        \toprule
        \textbf{When} & \textbf{What} \\
        \midrule
        Lecture & Parts 2, 3: build; the instructor visits each team \\
        Lab, Part A & Build, measure (80 min) \\
        Lab, Part B & Demonstrations (70 min) \\
        Lab, Part C & Report, feedback (30 min) \\
        \bottomrule
      \end{tabular}
    \end{column}
    \begin{column}{0.46\textwidth}
      \begin{block}{A demonstration of 10 minutes}
        \small
        \begin{itemize}
          \item The problem and the requirements
          \item The system works: live
          \item The benchmark table: latency, memory, energy
          \item Two decisions and their numbers
          \item Questions
        \end{itemize}
      \end{block}
    \end{column}
  \end{columns}
  \vskip 0.6em
  \small
  \begin{itemize}
    \item Finish the measurements before the polish. A working system with
          numbers scores more than a nice system with no numbers.
    \item Plan B: a recorded input or a video of the system, if the live
          demonstration fails.
  \end{itemize}
\end{frame}

% ------------------------------------------------------------------ exercise
\exerciseframe{One trade-off of your capstone}{%
  For your capstone system:
  \begin{enumerate}
    \item Name one trade-off of your design: what you gain and what you
          pay.
    \item Name the number that decides it, the unit, and how you measure
          it (which board, which window, how many runs or samples).
    \item Write the value that would change your decision.
  \end{enumerate}
  \vskip 0.4em
  Use the form of the Decision Log. You have 10 minutes. Two teams read
  their answer to the class.}

\answerframe{One trade-off of your capstone}{%
  \small
  Example answer for the machine monitor of Day 14:
  \begin{itemize}
    \item \textbf{Trade-off}: a battery on the machine (no cable) against
          the Wi-Fi link that the dashboard needs.
    \item \textbf{Number}: the mean current of the XIAO, measured with a
          USB power meter for 10 minutes in each mode. The budget says at
          least 167.2 mW with Wi-Fi connected (at most 22 h on 1000 mAh) and
          24.2 mW with light sleep and Wi-Fi each hour (6.4 days).
    \item \textbf{Change}: if the measured mean is below 5.1 mW, a battery
          of 1000 mAh lasts 30 days and the cable is not necessary.
  \end{itemize}
  A good answer names one number, its unit, its method, and a limit.}
